% Lösungen Einheit 3 von 5 – Dreieck (zusammenbau v0.1)
\einheitenkopf{Einheit 3 von 5 · Dreieck}

%% TODO Lösung der Erklärzeile 18g) fehlt in der Bank
\erg{18}{a) Strecke von C senkrecht auf AB \quad b) $A = 8 \cdot 7 : 2 = 28$ cm² \quad c) $A = 14 \cdot 9 : 2 = 63$ m² \quad d) $A = 6 \cdot 11 : 2 = 33$ dm² \quad e) $A = 20 \cdot 15 : 2 = 150$ cm² \quad f) $A = 9 \cdot 5 : 2 = 22{,}5$ m² \quad g) \ldots \quad h) $A = 15 \cdot 8 : 2 = 60$ cm² \quad i) $A = 6 \cdot 4 : 2 = 12$ cm² \quad j) $A = 4{,}6 \cdot 3 : 2 = 6{,}9$ cm² \quad k) $g = 2 \cdot 26 : 4 = 13$ cm \quad l) $u = 3s$ \quad m) $QC = 2 \cdot 120 : 16 = 15$ m}
\erg{19}{a) $A = 7{,}5 \cdot 6{,}4 : 2 = 24$ cm²}
\erg{20}{a) $u = 5 + 7 + 9 = 21$ cm}
\erg{21}{a) Jana hat vergessen zu halbieren. Richtig: $A = 9 \cdot 6 : 2 = 27$ cm² \quad b) Zwei gleiche Dreiecke lassen sich zu einem Parallelogramm mit derselben Grundseite und derselben Höhe zusammenlegen; ein Dreieck ist die Hälfte davon. \quad c) $A = 8{,}4 \cdot 3{,}5 : 2 = 14{,}7$ m²; $14{,}7 : 6 = 2{,}45$, also $3$ Dosen \quad d) gleichseitiges Dreieck mit drei Seiten $k$}

21d)

\dreieck{(0,0)}{(3,0)}{(1.5,2.6)}{$k$}{$k$}{$k$}{}{}{}

